\section{Analogy benchmarks: complete proofs}
\label{sec:annex-p19-analogy-and-problem-transfer}

These proofs concern the explicitly defined handbook models
in \cref{sec:p19-analogy-benchmarks}. They are not proofs
of the empirical transfer rates or of a psychological
processing theory.

\begin{annexitem}{annex:p19.analogy-structural-map}{thm:p19-analogy-structural-map}{The recursive substitution and its consistency}
\begin{proof}
Give each object expression height zero. Give a
constructor expression one plus the largest height
of its argument expressions; a constructor of arity
zero has height one. Every expression has finite
height by the definition of the free expression
language.

Define \(\tau_m(b)=m(b)\) for \(b\in B\).
If \(F(e_1,\ldots,e_r)\) is well formed and
substitution has been defined on its arguments,
put
\[
 \tau_m(F(e_1,\ldots,e_r))
          =F(\tau_m(e_1),\ldots,\tau_m(e_r)).
\]
For \(r=0\) this fixes the same nullary constructor.
Induction on height defines the substitution on
every expression. At each step the argument sorts
are unchanged by the inductive hypothesis and
the constructor is unchanged, so its required
sorts and resulting proposition sort are respected.
The construction preserves its arity and the
position of every argument. It also preserves
height and the entire tree shape.

Any substitution with the stated requirements
must have these same values on objects and
then on every constructor expression.
Induction therefore proves uniqueness.
In particular every occurrence of one object
\(b\) has the same image \(m(b)\), whatever
the surrounding constructor or its argument
position. The two occurrences of \(b\) in
the connected example must both become \(y\).
They cannot become different objects under
one function \(m\).

Suppose now that \(m\) is injective and that
\(\tau_m(e)=\tau_m(f)\).
An object image is an object and cannot equal
a constructor expression. If both are objects,
injectivity of \(m\) gives \(e=f\).
If both are constructor expressions, equality
in free syntax forces the same outer constructor,
the same arity and equal corresponding argument
images. Apply induction to those arguments
to conclude equality of all the original
arguments and hence of \(e,f\).
This proves injectivity on all expressions.

The displayed substitution of the two relations
and their higher-order link now follows by
three applications of the recursive rule.
The expression \(A(b)\) is a separate member
of the supplied representation, not a subexpression
of that link. Omitting it therefore leaves
the connected expression and its image unchanged.
This is a statement about this explicit syntax,
before any quotient by logical equivalence
or any psychological interpretation.
\end{proof}
\end{annexitem}

\begin{annexitem}{annex:p19.analogy-threshold-transfer}{thm:p19-analogy-threshold-transfer}{Exact resource feasibility and two failed transfers}
\begin{proof}
If \(0\le x_j\le h\) for every route, then
\[
                       \sum_{j=1}^k x_j\le kh.
\]
Thus success requires \(kh\ge q\).
Conversely, if \(kh\ge q\), assign
\(x_j=q/k\) to every route. These quantities
are nonnegative, at most \(h\), and sum
to \(q\); hence they form a feasible plan.

Since \(h>0\), the condition on the positive
integer \(k\) is \(k\ge q/h\). Its least
solution is the ceiling \(\lceil q/h\rceil\).
The assumption \(q>h\) makes this at least
two. If \(\lambda>0\), multiplication of
every amount and both thresholds by
\(\lambda\) preserves each inequality,
and division by \(\lambda\) proves the
converse. This is the asserted scaling
equivalence.

For \(h=4,q=10,k=3\), each of the three
amounts \(10/3\) is less than \(4\), and
their sum is exactly \(10\). A ten-unit
load on a single route violates its
four-unit capacity.
With two routes the total safely deliverable
amount is at most eight, so the target
requirement cannot be reached. With three
routes but a maximum rather than an
additive aggregation rule, the effect
is at most
\(\max_jx_j\le h<q\), here \(4<10\). No choice of safe
amounts can then succeed.

These two targets can retain the same
description of a central goal and resources
approaching it while changing a premise
essential to the plan's proof. The successful
source allocation therefore transfers only
when the route and aggregation conditions
used in that proof are also established
for the target.
\end{proof}
\end{annexitem}

\begin{annexitem}{annex:p19.analogy-candidate-truth}{thm:p19-analogy-candidate-truth}{Two agreeing target descriptions with opposite candidate truth}
\begin{proof}
Use the three-element domain \(\{x,y,z\}\),
with each displayed name denoting its
corresponding element. In both structures
interpret
\[
                  R=\{(x,y)\},\qquad S=\{(y,z)\}.
\]
In \(\mathcal M_0\) set \(P=\varnothing\).
In \(\mathcal M_1\) set \(P=\{(x,z)\}\).
Both structures satisfy all known target
facts. The same external object correspondence
\(m(a)=x,m(b)=y,m(c)=z\) can be used with
either. Recursive substitution produces
\(\tau_m(P(a,c))=P(x,z)\) in both cases.
Yet \(\mathcal M_0\models\neg P(x,z)\) and
\(\mathcal M_1\models P(x,z)\).
Consequently neither candidate truth value
is a semantic consequence of the known
target facts.

A higher-order expression is also produced
by a rule of syntax. In the stated expression
language no causal or explanatory truth
clause has been imposed on
\(\mathsf{Supports}\).
For instance, even after giving both
argument propositions truth value true,
one can interpret that constructor as a
Boolean operation constantly false or
as one constantly true; its expression
tree is identical in the two interpretations.
A substantive interpretation may impose
further conditions, but those conditions
are additional premises. The structural
substitution itself supplies none.
\end{proof}
\end{annexitem}
